\secciongrande{Modelos de ejecución de código: Los cuatro campos del diseño web}

\resumen{
  El desarrollo web moderno se fundamenta en el modelo cliente-servidor y en cuatro campos interdependientes del diseño: el visual, el tecnológico, el de distribución y el de propósito. Desde los orígenes del HTML estático hasta las aplicaciones dinámicas actuales, el programador debe dominar un campo multidisciplinar que abarca desde la captación de necesidades hasta la presentación final, garantizando estándares mantenidos por el W3C.
}

\subtituloA{Origen y alcance de la web}

En sus orígenes, la creación de una página web requería únicamente el lenguaje HTML, cuyas marcas permitían construir documentos sencillos que cualquier navegador podía visualizar. Estas páginas se alojaban en servidores conectados permanentemente, transmitiéndose a los equipos clientes que las solicitaban.

La web se fundamenta en el concepto de hipertexto: palabras o frases que, al ser seleccionadas, conducen a información relacionada. Ante el auge de esta tecnología, se creó el consorcio W3C, organización encargada de desarrollar y mantener los estándares web que garantizan la interoperabilidad entre diferentes plataformas y navegadores.

Un programador web debe dominar todas las fases de un proyecto: captar las necesidades del cliente, definir objetivos y estrategia, desarrollar la interfaz gráfica, elaborar prototipos y realizar la presentación final. Es imprescindible conocer el orden secuencial de cada fase, ya que el desarrollo web constituye un campo inherentemente multidisciplinar.
